% LSPV-1 REPORT::QNORM1::2026-10-12T20:08:44Z
% !TEX program = xelatex
\documentclass[11pt]{article}
\usepackage[a4paper,margin=0.88in]{geometry}
\usepackage{fontspec}
\usepackage{amsmath,amssymb,amsthm,booktabs,longtable,array,xurl,hyperref,microtype}
\hypersetup{colorlinks=true,linkcolor=blue,urlcolor=blue,citecolor=blue}
\setlength{\emergencystretch}{3em}
\newtheorem{theorem}{Theorem}
\newtheorem{corollary}{Corollary}
\newcommand{\SAT}{\mathsf{SAT}}
\newcommand{\Pclass}{\mathsf P}
\newcommand{\NPclass}{\mathsf{NP}}
\newcommand{\EXACT}{\textbf{PROVED/EXACT}}
\newcommand{\KNOWN}{\textbf{KNOWN THEOREM}}
\newcommand{\COMP}{\textbf{COMPUTATIONAL}}
\newcommand{\COND}{\textbf{CONDITIONAL}}
\newcommand{\CAND}{\textbf{CANDIDATE}}
\newcommand{\OPEN}{\textbf{OPEN}}
\newcommand{\RET}{\textbf{FALSIFIED/RETIRED}}
\title{QNORM1: Exact Quantifier Normal Form for a Universal Clean-Core Time Lower Bound}
\author{LSPV-1 --- Universal Core Lower Bound}
\date{13 October 2026 (KST)}

\begin{document}
\maketitle

\begin{center}
\fbox{\parbox{0.94\linewidth}{
\textbf{Executive status.} P versus NP remains \OPEN.  This run resolves the
newly named \texttt{TOTAL-COST-CLEAN-CORE-OBSTRUCTION1} when ``total cost''
dominates ordinary running time.  On a fixed polynomial-time many-one
NP-complete clean core, the fully quantified cofinal superpolynomial-time
statement is exactly equivalent to \(P\ne NP\).  This is a quantifier
normalization and reconciliation of the earlier COREEQUIV1 theorem, not a new
separation.  The target is legitimate as a final endpoint but is \RET{} as a
weaker intermediate bridge.  Fixed-exponent, per-length-machine, and
phase-local variants do not recover the endpoint.  A bounded 28,475-check
audit passed with zero failures and is only \COMP.}}
\end{center}

\section{Predecessor write-recovery preflight}

LSPV-1's predecessor is LSPV-5.  The newest verified canonical MAIN readable
from GitHub is
\path{research/main/2026-10-13_0345_MAIN_REPPROMISE1.tex}.  The immediately
preceding synthesis invocation produced no newer verified canonical MAIN.
The exact later scientific body was unavailable, so it was not reconstructed.
The nonblocking recovery request
\path{research/review/outbox/lower_bound/2026-10-13_0508_LSPV5_WRITE_RECOVERY.md}
was written and verified at commit
\texttt{27ad398d61052d51f9d4fdf6e44da1347ed77621}.  Evidence supports
\textbf{PROMPT / PROCEDURE INTERRUPTION or UNKNOWN}, not a claim about the
missing science.  Current research continued as required.

\section{Exact family, encoding, machine, and target}

Fix a decidable syntactic domain \(D_K\subseteq\{0,1\}^*\) and a language
\(K\subseteq D_K\).  Inputs use one fixed complete self-delimiting binary
encoding and \(N=|x|\).  Assume a polynomial-time many-one reduction
\(r:\{0,1\}^*\to D_K\) and a polynomial \(q\) satisfy
\[
 z\in\SAT\iff r(z)\in K,
 \qquad |r(z)|\le q(|z|).                                  \tag{1}
\]
The project instance is the duplicate-free signed Monotone
\(3\)SAT-\((2,2)\) core, whose restricted satisfiability problem is
NP-complete \cite{Docker19}.

Let \(\operatorname{Dec}(K)\) contain total uniform deterministic ordinary
machines correct on every \(x\in D_K\).  Input reading, preprocessing,
simulation, indexing, extraction, replay, reconstruction, and the one-bit
output are charged in \(T_A(x)\).  Define
\begin{align}
 \mathsf{StrongWC}(K):\quad
 &\forall A\in\operatorname{Dec}(K)\;\forall d\ge1\;
   \forall C\ge1\;\forall N_0\;\exists x\in D_K:\nonumber\\[-2mm]
 &\hspace{37mm}|x|\ge N_0\ \land\ T_A(x)>C|x|^d,          \tag{2}\\
 \mathsf{UnitWC}(K):\quad
 &\forall A\in\operatorname{Dec}(K)\;\forall d\ge1\;
   \forall N_0\;\exists x\in D_K:\nonumber\\[-2mm]
 &\hspace{37mm}|x|\ge N_0\ \land\ T_A(x)>|x|^d.          \tag{3}
\end{align}

\section{Exact normal-form theorem}

\begin{theorem}[\texttt{CLEAN-CORE-TIME-QUANTIFIER-NORMAL-FORM1}; \EXACT]
Under (1), the following are equivalent:
\[
 P\ne NP,
 \qquad K\notin\Pclass,
 \qquad \mathsf{StrongWC}(K),
 \qquad \mathsf{UnitWC}(K).                               \tag{4}
\]
\end{theorem}

\begin{proof}
Because \(K\in NP\) and (1) makes it NP-hard, \(K\notin P\) is equivalent to
\(P\ne NP\).  Clearly (2) implies (3) by setting \(C=1\).

If (3) fails, quantifier negation gives one total uniform exact decider
\(A\), one \(d\), and one \(N_0\) such that
\[
 \forall x\in D_K\quad |x|\ge N_0\Longrightarrow T_A(x)\le |x|^d. \tag{5}
\]
There are finitely many shorter strings.  Totality makes the maximum running
time of \(A\) on their intersection with \(D_K\) a finite constant.  Thus
\(A\) decides \(K\) in polynomial time, so \(K\in P\).

Conversely, suppose (2) fails.  Then some \(A,d,C,N_0\) obey
\[
 \forall x\in D_K\quad |x|\ge N_0\Longrightarrow
 T_A(x)\le C|x|^d.                                         \tag{6}
\]
The same finite-exception maximum makes \(A\) polynomial on all of \(D_K\).
Composing the actual reduction \(r\) with this same actual machine \(A\)
decides SAT in time polynomial in \(|z|\), since reduction time and
\(Cq(|z|)^d\) are both polynomial.  Hence \(P=NP\).  Taking contrapositives
proves \(K\notin P\Rightarrow(2)\), completing all equivalences.
\end{proof}

\begin{corollary}[Counterfactual under \(P=NP\); \EXACT]
If \(P=NP\), then for this fixed clean core there exist one uniform exact
decider \(A\), constants \(d,C,N_0\), and the fully charged bound
\[
 \forall x\in D_K,quad |x|\ge N_0\Longrightarrow
 T_A(x)\le C|x|^d.                                         \tag{7}
\]
Thus \(P=NP\) implies the \emph{negation} of the universal obstruction.  No
solver-independent implication from \(P=NP\) to (2) can be used as an
intermediate lemma without already deriving the desired contradiction.
\end{corollary}

The result is the exact constant-explicit normal form of COREEQUIV1.  The
earlier theorem already proved the unit-coefficient form (3), including the
finite exceptional set and reduction cost.  QNORM1's new contribution is the
explicit equivalence with (2) and the classification of the newly named
whole-execution target; it does not relabel COREEQUIV1 as new.

\section{Falsified weakenings and exact repairs}

\begin{longtable}{p{0.31\linewidth}p{0.19\linewidth}p{0.42\linewidth}}
\toprule Variant & Status & Exact defect or repair\\ \midrule
All \(A,d,C,N_0\), one later witness & \EXACT & This is (2), equivalent to
the separation on the fixed NP-complete core.\\
One fixed exponent \(c\) & \RET{} as P-vs-NP endpoint & Excluding
\(O(N^c)\) does not exclude an \(O(N^{c+1})\) decider.  It separates only
from the named fixed class unless another theorem supplies exponent collapse.\\
\(\forall N\,\exists A_N\) & \RET & A separate finite hardwired machine can
answer each length.  This is POINTINF1's uniformity failure, not one program.\\
No cofinal threshold and no multiplicative constant & \RET & A finite
small-input timing anomaly can witness every tested exponent; restore
\(\forall C\forall N_0\).\\
One common hard sequence for all programs & \CAND{} at best & Recognizable
YES/NO sequences admit shortcut wrappers.  The correct endpoint lets the
hard instance depend on \(A,d,C,N_0\).\\
Phase-local work after complete source coverage & \RET & PHASEBRIDGE1's
cofactor compiler moves decision work before the final reads.\\
Arbitrary representation or padding & \RET & Fix the encoding, or prove
polynomially bounded translations and charge them.\\
\bottomrule
\end{longtable}

The strict order in (2) is important:
\[
 \forall A\;\forall d\;\forall C\;\forall N_0\;\exists x.
\]
Moving \(A\) after the length, fixing \(d\), or demanding one solver-independent
hard sequence changes the theorem.

\section{Semantic compatibility and SAME\_REALIZATION}

For a formula \(F\) and partial assignment \(\alpha\), define the external
compatibility predicate
\[
 H(F,\alpha)={\bf 1}[F\upharpoonright\alpha\text{ is satisfiable}]. \tag{8}
\]
Nothing in the proof asserts that an arbitrary decider evaluates all values
of \(H\), exposes a witness, materializes a graph, preserves a cut state, or
communicates a transcript.  The output contract is one bit.  Semantic-state
count is never converted into runtime.

SAME\_REALIZATION is explicit only where it is justified: in the reverse
direction, the composed SAT program contains the reduction, serialization,
the actual run of the same \(A\), and the finite exceptional cases.  External
extraction of \(H\) would be a separate charged computation and cannot be
attributed to \(A\)'s trace.  PHASEBRIDGE1 therefore remains compatible with
(4): it refutes a phase-local exposure claim but not the whole-execution final
endpoint.

\section{Bounded reproducible hostile audit}

The dependency-free command
\begin{verbatim}
timeout 30s python3 \
 experiments/review/clean_core_time_quantifier_normal_form_audit.py \
 --output results/raw/review/2026-10-13_clean_core_time_quantifier_normal_form.json
\end{verbatim}
completed in \(0.00578\) seconds as one chunk.  It performed 28,475 checks
with zero failures.  Controls were 2SAT, Horn, affine Gaussian elimination,
bounded treewidth, and independent products under DIRECT, COPY, DEAD, DELAY,
INTERP, and RECOMP wrappers.  It also checked fixed-exponent polynomial
countermodels, finite-threshold anomalies, pointwise machine selection,
polynomial encoding conversions, and bounded superpolynomial proxies.

This audit is \COMP{} only.  The symbolic proof above, not finite data,
establishes (4).  Script SHA-256:
\nolinkurl{5e45c6c0b638a840ec68f3ccfd9dbdb6196f5c52829a34e4f3cccde31c1e42ac}.
JSON SHA-256:
\nolinkurl{a2d0d66ce829dc94807b45513ffa22703dfe82912b92e31a2de8f0a95bc33bbf}.

\section{Primary-literature and barrier-transfer preflight}

The first obstacle is logical strength: a whole-execution superpolynomial
lower bound for an NP-complete language is the separation, not a structural
precursor.  Cook's NP-completeness theorem and Karp's reduction architecture
are directly transferable to the reduction step \cite{Cook71,Karp72}.  The
second obstacle is obtaining any weaker model-matched unconditional endpoint
for arbitrary SAT computation.  FLMV gives an unconditional simultaneous
time--space exclusion only in its stated random-access model and parameter
region \cite{FLMV05}; a compiler landing the same realization there remains
missing.

\begin{longtable}{p{0.28\linewidth}p{0.21\linewidth}p{0.42\linewidth}}
\toprule Source/domain & Classification & Correspondence and missing theorem\\ \midrule
Cook--Karp reductions & directly transferable theorem & Equation (1) plus
polynomial output length transfers a polynomial core decider to SAT.\\
FLMV time--space lower bound & useful theorem requiring a new transfer & It
constrains one cited random-access realization; arbitrary clean-core deciders
are not automatically in its forbidden region.\\
Time hierarchy & useful structural analogy requiring a new transfer & A
strict same-model speedup is required; HIERFIRE1 shows same-scale
materialization retains the unknown solver exponent.\\
Communication, VLSI, branching programs, proof complexity & useful
structural analogies requiring a new transfer & Cuts, routing, widths, and
proof sizes constrain restricted objects, not every ordinary one-bit decider.\\
Complexity cores assuming \(K\notin P\) & misleading analogy / failed
transfer & For this NP-complete \(K\), the premise already assumes the
desired separation.\\
\bottomrule
\end{longtable}

The proof of (4) relativizes whenever the NP-hard reduction and decider are
relativized together, so it does not evade Baker--Gill--Solovay
\cite{BGS75}.  It also algebrizes under the analogous extension and supplies
no nonalgebrizing hinge \cite{AW09}.  It asserts no large constructive circuit
property, so Natural Proofs is not directly invoked; a future unrestricted
circuit endpoint would require a separate audit.

\section{Post-quantum technique ledger}

Schur squares, puncturing/shortening, quadratic-relation rank spectra,
Hilbert functions, syzygies, and graded Betti numbers are
\textbf{FAILED TRANSFERS} to (2): no theorem forces an arbitrary decider to
expose the corresponding code or algebra, and hidden algebraic structure may
be an algorithmic weakness.  Collision entropy is an
\textbf{EXACT OBJECT/THEOREM TRANSFER} for a declared map
\(L:A\to F_{\rm acc}\), including
\[
 |\operatorname{im}L|\ge |A|^2/
 |\{(y,y'):L(y)=L(y')\}|,
\]
but it does not lower-bound the runtime of an arbitrary program.  LWE-style
black-box extraction and special-soundness/collision extraction are
\textbf{ANALOGY/CANDIDATE needing a new bridge}; extraction, queries, replay,
and reconstruction must be charged.  Cryptographic binding, hashes, and
random oracles are \textbf{FAILED TRANSFERS} as unconditional endpoints.

\section{CORPUS REVIEWED}

The current canonical GitHub tree was surveyed.  The run read root
\path{main.tex}, \path{AGENTS.md}, \path{CODEX.md}, the complete PNP research
skill, canonical \path{research/review/PROTOCOL.md}, actual \path{main4.tex}
and \path{main5.tex}, and all TASK-001--TASK-007 specifications.  It read the
newest verified MAIN REPPROMISE1; newest verified reports HIERFIRE1,
COVERCOF1, LOGCOF1, and PHASEBRIDGE1; and materially relevant historical
COREEQUIV1, JOINTENDPOINT1, CIRCEQUIV1, POINTINF1, and SYNTHESIS-0022.

Repository searches covered quantifier, padding, wrapper, easy-island,
complexity-core, circuit, hierarchy, and whole-execution formulations.  The
HEC, VSC-HEC, G-PPC-HEC, CRS-PPC, CCSPC, RCCN, and Search-G causal branch,
governance notes, and cited experiment/result paths were surveyed for the
arbitrary-algorithm bridge.  Optional skill references
\path{references/modes.md} and
\path{references/research-protocol.md} were unavailable locally; the complete
skill, AGENTS, and canonical protocol governed instead.

\section{DUPLICATION / CONTRADICTION AUDIT}

COREEQUIV1 already proves (3) equivalent to \(P\ne NP\).  JOINTENDPOINT1
already proves that any fixed joint gauge with a positive time power gives the
same endpoint.  CIRCEQUIV1 classifies an unrestricted circuit endpoint as the
stronger \(NP\not\subseteq P/poly\), and POINTINF1 rejects per-length machine
selection.  These results are adopted, not duplicated.

PHASEBRIDGE1's newly named whole-execution target is not a contradiction: its
phase-local compatibility claim is retired, while its remaining direct time
claim is precisely (3).  Root MAIN already states that cofinal
superpolynomial time on the corrected NP-hard core is separation-equivalent.
The only new mathematical increment is the constant-explicit normal form (2)
and the associated quantifier-repair table.

\section{WHAT CHANGED RELATIVE TO ROOT MAIN}

The pre-run LSPV-1 region correctly labels the direct clean-core time target
OPEN and calls joint superpolynomial time separation-strength.  This run adds
the exact \(C\)-explicit quantifiers, identifies
\texttt{TOTAL-COST-CLEAN-CORE-OBSTRUCTION1} with that existing endpoint when
total cost dominates time, and retires its presentation as a weaker
intermediate.  The direct statement itself remains \OPEN{} because proving it
would prove \(P\ne NP\).  No exponent, family, or machine model is silently
strengthened.

Actual \path{main4.tex} remains the ARC achievable-region manuscript and
actual \path{main5.tex} remains the WRAPPER construction manuscript.  Their
partial promise regions and recovery rules neither prove nor refute (2).

\section{ADOPTED / REJECTED / RETIRED PRIOR COMPONENTS}

\textbf{Adopted:} the corrected signed NP-complete core; COREEQUIV1;
JOINTENDPOINT1; POINTINF1's uniformity guard; PHASEBRIDGE1's cofactor bypass;
HIERFIRE1's exponent ledger; and FLMV only in its exact model.

\textbf{Rejected as transfers:} internal compatibility exposure, semantic
state counts, proof/circuit/communication size without a compiler, and a
complexity-core theorem whose premise already assumes \(K\notin P\).

\textbf{FALSIFIED/RETIRED:}
\texttt{TOTAL-COST-CLEAN-CORE-OBSTRUCTION-AS-INTERMEDIATE1}; fixed-exponent
and per-length-machine substitutes; and phase-local post-coverage work.  The
fully quantified direct endpoint (2) is not retired; it remains \OPEN.

\section{STALE OR OVERSTATED CLAIMS FOUND}

Calling the whole-execution superpolynomial-time target merely a ``next
bridge'' is overstated: under the fixed NP-complete encoding it is the final
separation.  Calling a fixed-exponent lower bound sufficient is also
overstated because the hypothetical polynomial exponent is arbitrary.  Any
claim that source coverage forces paid compatibility contradicts the verified
cofactor compiler.  No current theorem forces arbitrary SAT computation to
expose \(H\), a witness, a transcript, or a named memory state.

\section{Strict status ledger and handoff}

\begin{center}\small
\begin{tabular}{p{0.57\linewidth}p{0.25\linewidth}}
\toprule Claim & Status\\ \midrule
CLEAN-CORE-TIME-QUANTIFIER-NORMAL-FORM1 & \EXACT\\
Signed clean-core NP-completeness & \KNOWN\\
28,475-check hostile audit & \COMP\\
TOTAL-COST obstruction as weaker intermediate & \RET\\
Fully quantified direct endpoint (2) & \OPEN\\
Same-realization landing in a cited restricted lower bound & \OPEN\\
Nonrelativizing/nonalgebrizing hinge & \OPEN\\
\(P\ne NP\) & \OPEN\\
\bottomrule
\end{tabular}
\end{center}

The highest-value next theorem is not another restatement of (2).  It is a
genuinely model-matched implication that, from an arbitrary polynomial SAT
decider with unknown exponent, constructs one fully charged ordinary
realization inside a cited unconditional forbidden region.  The same output
machine must bear source reads, materialization, failed branches, queries,
replay, reconstruction, output, time, and peak workspace.  This load-bearing
arrow remains \OPEN.

\begin{thebibliography}{9}
\bibitem{Cook71}
S. A. Cook, ``The Complexity of Theorem-Proving Procedures,''
\emph{STOC 1971}, 151--158. doi:10.1145/800157.805047.

\bibitem{Karp72}
R. M. Karp, ``Reducibility Among Combinatorial Problems,'' in
\emph{Complexity of Computer Computations}, 1972, 85--103.

\bibitem{Docker19}
J. D\"ocker, ``Monotone 3-Sat-(2,2) is NP-complete,'' 2019.
\url{https://arxiv.org/abs/1912.08032}.

\bibitem{FLMV05}
L. Fortnow, R. Lipton, D. van Melkebeek, and A. Viglas,
``Time-Space Lower Bounds for Satisfiability,'' \emph{JACM} 52(6),
835--865, 2005. doi:10.1145/1101821.1101822.

\bibitem{BGS75}
T. Baker, J. Gill, and R. Solovay, ``Relativizations of the \(P=?NP\)
Question,'' \emph{SIAM J. Comput.} 4(4), 431--442, 1975.
doi:10.1137/0204037.

\bibitem{AW09}
S. Aaronson and A. Wigderson, ``Algebrization: A New Barrier in Complexity
Theory,'' \emph{ACM TOCT} 1(1), 2009. doi:10.1145/1490270.1490272.
\end{thebibliography}

\bigskip
\noindent\textbf{Final status: P versus NP remains OPEN.}

% END_OF_VERIFIED_REPORT::LSPV-1::QNORM1
\end{document}
